\chapter{Nuclear Synthesis in Stars}

\section{Nuclear Basics}

\subsection{Nuclear Size and Binding Energy}
\begin{theory}
\begin{align}
    R &= R_0A^{1/3}, \qquad R_0\approx1.2\times10^{-15}\,\text{m}, \\
    BE(Z,N) &= \left[Zm_p+Nm_n-M(Z,N)\right]c^2.
\end{align}
Binding energy per nucleon increases for light nuclei and decreases for very heavy nuclei due to Coulomb repulsion.
\end{theory}

\subsection{Why Fusion Powers Stars}
\begin{intuition}
Stars gain energy when fusing light nuclei toward higher binding energy per nucleon.
This trend reverses near iron, where further fusion is endothermic.
\end{intuition}

\section{Hydrogen and Helium Burning}

\subsection{Proton--Proton Chain}
\begin{theory}
Representative net reaction:
\begin{equation}
    4\,{}^1\text{H}\rightarrow{}^4\text{He}+2e^++2\nu_e+\text{energy}.
\end{equation}
Total released energy is about $26.7\,$MeV, with part carried away by neutrinos.
\end{theory}

\subsection{CNO Cycle}
\begin{theory}
CNO uses C, N, O as catalysts for H-burning and is much more temperature sensitive than pp burning.
Near solar-core temperatures, pp-chain is weaker in $T$-sensitivity ($\sim T^4$) while CNO can scale roughly as $T^{\sim 18}$.
\end{theory}

\subsection{Triple-Alpha Process}
\begin{theory}
\begin{equation}
    3\,{}^4\text{He}\rightarrow{}^{12}\text{C}+\gamma,
\end{equation}
enabled by resonance through unstable ${}^8\text{Be}$ and an excited Hoyle state in ${}^{12}\text{C}$.
Typical ignition requires $T\gtrsim10^8\,$K.
\end{theory}

\begin{importantnote}
The first pp step ($p+p\to d+e^++\nu_e$) is a weak-interaction bottleneck.
It is slow enough to set long stellar main-sequence lifetimes.
\end{importantnote}

\section{Barrier Penetration and Reaction Rates}

\subsection{Coulomb Barrier and Tunneling}
\begin{foundation}
Classically, nuclei with $E<V_C$ cannot fuse.
Quantum tunneling gives non-zero penetration probability,
\begin{equation}
    P\sim \exp\left[-\sqrt{\frac{E_G}{E}}\right],
\end{equation}
where $E_G$ is the Gamow energy.
\end{foundation}

\subsection{Thermally Averaged Cross Section}
\begin{theory}
\begin{equation}
    \langle\sigma v\rangle = \left(\frac{8}{\pi\mu}\right)^{1/2}
    \frac{1}{(k_BT)^{3/2}}\int_0^\infty S(E)
    \exp\left[-\frac{E}{k_BT}-\sqrt{\frac{E_G}{E}}\right]dE.
\end{equation}
The integrand peaks at the Gamow energy window, balancing high-energy tail abundance and tunneling probability.
\end{theory}

\subsection{Temperature Sensitivity}
\begin{theory}
Near a reference temperature, reaction rates are often approximated by
\begin{equation}
    R\propto T^{\nu},
\end{equation}
with small $\nu$ for pp-chain and much larger $\nu$ for CNO and triple-alpha.
This sensitivity drives core structure differences between low-mass and massive stars.
\end{theory}

\begin{example}{Qualitative core response to heating}{ex:temp-feedback}
Suppose core temperature rises slightly.

\smallskip\textbf{Discussion.}
In a pp-dominated core, energy generation increases moderately, giving mild feedback.
In a CNO-dominated core, generation rises steeply, making burning more centrally concentrated and favoring convective cores.
\end{example}

\section{Advanced Burning and Stellar Endpoints}

\subsection{Late Burning Stages in Massive Stars}
\begin{theory}
After He burning, successive C, Ne, O, and Si burning stages occur at progressively higher temperatures and shorter timescales.
Fusion energy production terminates near iron-group nuclei.
\end{theory}

\subsection{Neutrino Cooling and Rapid Evolution}
\begin{extra}
At high temperature, neutrino processes (e.g., pair production) carry energy away efficiently.
This accelerates late-stage evolution and contributes to core collapse in massive stars.
\end{extra}

\section{Solar Neutrino Problem}
\begin{theory}
Historically, measured solar neutrino flux was below early predictions.
The resolution is flavor oscillation:
\begin{equation}
    \nu_e\leftrightarrow\nu_\mu\leftrightarrow\nu_\tau,
\end{equation}
so detectors sensitive mostly to $\nu_e$ observed an apparent deficit.
\end{theory}

\begin{remark}
Solar neutrino measurements are both a test of stellar fusion models and direct evidence for neutrino mass and mixing.
\end{remark}
